%% RELIC — Conference Paper (Submission-Ready Draft)
%% Target: ICSE / ESEC/FSE / ASE / ICSME  (10-page double-column ACM/IEEE format)
%% Compilation: pdflatex paper.tex (requires booktabs, amsmath, graphicx, xspace)
%%
%% DOUBLE-BLIND SUBMISSION: author and affiliation fields are left as placeholders.
%% Remove "Anonymous" markers and insert real names before camera-ready.
%%
%% All numbers in this paper are traced to results/FINAL_RESULTS_MANIFEST.json (v1.4).
%% Authoritative source: make reproduce-paper1

\documentclass[sigconf,anonymous,review]{acmart}

%% Packages
\usepackage{booktabs}
\usepackage{amsmath}
\usepackage{amssymb}
\usepackage{graphicx}
\usepackage{xspace}
\usepackage{microtype}
\usepackage{hyperref}
\usepackage{cleveref}
\usepackage{listings}
\usepackage{xcolor}
\usepackage{multirow}

%% Convenience macros
\newcommand{\bcg}{\textsc{BCG}\xspace}
\newcommand{\system}{\textsc{RELIC}\xspace}
\newcommand{\carddemo}{\textit{CardDemo}\xspace}
\newcommand{\gnucobol}{\textit{GnuCOBOL}\xspace}
\newcommand{\notobs}{\texttt{NOT\_OBSERVED}\xspace}
\newcommand{\supported}{\texttt{STATIC\_DYNAMIC\_SUPPORTED}\xspace}
\newcommand{\conflict}{\texttt{STATIC\_DYNAMIC\_CONFLICT}\xspace}
\newcommand{\ambiguous}{\texttt{AMBIGUOUS}\xspace}
\newcommand{\partiallysupported}{\texttt{PARTIALLY\_SUPPORTED}\xspace}
\newcommand{\ie}{\textit{i.e.},\xspace}
\newcommand{\eg}{\textit{e.g.},\xspace}

%% ACM metadata — anonymized for double-blind review
\acmConference[Submission]{Under Review}{2025}{}
\acmBooktitle{Manuscript under review}

\begin{document}

%%=======================================================================
\title{RELIC: Execution-Grounded Behavioral Analysis\\for Legacy Software Modernization}
%%=======================================================================

%% Anonymous submission — remove for camera-ready
\author{Anonymous Authors}
\affiliation{\institution{Anonymous Institution}}
\email{anonymous@anonymous.org}

\begin{abstract}
Legacy batch systems written in COBOL and Job Control Language (JCL) encode implicit
behavioral contracts across heterogeneous artifacts: temporal ordering obligations in
step sequences, resource availability guarantees in dataset disposition parameters,
and control-flow guards in return-code conditions.
These obligations are never explicitly represented, making semantic behavioral
regressions the dominant failure mode during modernization to cloud-native
orchestration platforms.

We present the \textbf{Behavioral Contract Graph} (\bcg), a provenance-aware formal
intermediate representation ($\mathit{BCG} = (V, E, C, P, D, \Gamma)$) that makes
implicit behavioral obligations explicit across JCL and COBOL artifacts.
Key to our approach is \emph{confidence-delta detection}: a contract can be
\emph{weakened} without being removed, and this weakening is a modernization risk
invisible to structural dependency analysis.

We evaluate \system on the AWS \carddemo mainframe benchmark across three research
questions.
\textbf{RQ1} (contract recovery): the TEMPORAL extraction pipeline achieves
micro-precision $= 1.00$, micro-recall $= 1.00$, micro-F1 $= 1.00$ against an
independently implemented deterministic oracle on 33 held-out JCL files
(56 TP, 0 FP, 0 FN; macro-F1 $= 0.58$, 95\% CI: [0.39, 0.73]).
This evaluation covers TEMPORAL contracts; RESOURCE and CONTROL\_FLOW require
human annotation.
\textbf{RQ2} (risk detection): against 104 evaluated fault-grounded mutations,
static \bcg achieves F1 $= 1.00$ and FPR $= 0.00$ on a 144-item
semantics-preserving negative benchmark.
The discriminating contribution is confidence-delta detection, which catches
8 \texttt{ALTER\_ERROR\_BRANCH} mutations invisible to both structural analysis
(Baseline~A, F1\,$=$\,0.90) and \bcg without confidence-delta (Baseline~E,
F1\,$=$\,0.96).
On this blinded benchmark, the evaluated LLM-only and Dependency+LLM baselines
achieved F1\,$=$\,0.038 and F1\,$=$\,0.109, respectively (lower bounds due to
97--98\% API error rates; 3 and 7 successful calls from 248 total).
\textbf{RQ3} (dynamic evidence): with 40 traces (2 real GnuCOBOL, 38 simulated
JCL step-ordering; \textbf{not z/OS}),
91.8\% of contracts reach \supported; the coverage experiment
empirically confirms $\notobs \neq \text{CONTRADICTED}$ across all coverage levels.
All FPR=0 results on the 144-item negative benchmark are consequences of the
closed-set benchmark design; open-world false-positive rates are not measured.
\end{abstract}

%%=======================================================================
\section{Introduction}\label{sec:intro}
%%=======================================================================

Mainframe batch systems underpin global financial infrastructure.
Over 220 billion lines of COBOL remain in active production, processing more than
95 billion transactions daily~\cite{comella2000survey}.
Modernizing these systems --- migrating JCL job streams to AWS Step Functions or
Apache Airflow, translating COBOL to Java, or re-hosting to cloud-native batch
engines --- is a high-value, high-risk engineering activity.
Industry postmortems consistently identify \emph{semantic behavioral regressions} as
the dominant failure mode: the modernized system runs without crashing but silently
violates an obligation that the original JCL enforced~\cite{ulrich2010information}.

\textbf{The root cause is structural.}
The behavioral obligations embedded in JCL artifacts --- step ordering, dataset
availability sequencing, return-code propagation chains --- are never made explicit.
When a migration engineer translates \texttt{COND=(4,GT,STEP01)} to an orchestration
condition, they must infer its semantics from context; there is no formal
specification to consult.
When an initialization step is dropped because it ``only creates an empty file,''
the implicit resource dependency between that step and the downstream transaction
programs is silently broken.

\textbf{The gap in existing approaches.}
Structural dependency analysis~\cite{tip1995slicing,reps1997profiling} identifies
\emph{what depends on what}, but cannot represent \emph{what behavioral obligations
must be preserved}.
Consider a modernization that removes the \texttt{COND} parameter controlling
step bypass on error (\texttt{ALTER\_ERROR\_BRANCH}).
The structural dependency graph is \emph{unchanged}: the same steps reference the
same datasets.
Structural Baseline A (program slicing) and BCG-without-confidence-delta
(Baseline E) both miss this class of regression.
Only explicit behavioral contract representation --- and confidence-delta detection
--- catches it.

\textbf{Our contribution.}
We introduce the \textbf{Behavioral Contract Graph (\bcg)}, a formal intermediate
representation that makes behavioral obligations explicit:
\begin{enumerate}
\item A typed entity graph over JCL jobs, steps, programs, and datasets;
\item Behavioral contracts (\textsc{Temporal}, \textsc{Resource},
      \textsc{Control-Flow}) with deterministic provenance;
\item A static/dynamic reconciliation mechanism that distinguishes
      \supported, \conflict, and \notobs;
\item A confidence-delta differencing algorithm (\texttt{BCGDiff}) that detects both
      \emph{broken} and \emph{weakened} contracts.
\end{enumerate}

\textbf{Research questions:}
\begin{description}
\item[\textbf{RQ1}] Can TEMPORAL behavioral contracts be automatically recovered
  with precision and recall sufficient for practical use, validated against an
  independently implemented oracle?
\item[\textbf{RQ2}] Does contract-aware BCG analysis improve modernization-risk
  detection compared with structural, complexity, LLM-based, and static-BCG
  baselines?
\item[\textbf{RQ3}] Does dynamic execution evidence increase contract validation
  coverage while preserving zero false positives on a semantics-preserving benchmark?
\end{description}

\textbf{Motivating example.}
\Cref{lst:motivating} shows an \texttt{ALTER\_ERROR\_BRANCH} mutation in JCL.
The original \texttt{COND=(4,GT,STEP01)} means: ``skip STEP02 if STEP01 returned a
code greater than 4.''
The modernized version changes the threshold to \texttt{COND=(16,GT,STEP01)},
allowing STEP02 to execute even when STEP01 signaled a soft error (RC=8--15).
The structural dependency graph is \emph{identical} in both versions --- STEP02
still follows STEP01.
Structural Baseline A misses this.
\bcg represents a \textsc{Control-Flow} contract with
$\mathit{static\_support} = 1.0$; the mutation reduces it to $0.7$.
BCGDiff detects the confidence-delta $\Delta = 0.3 > \theta$, raising a risk flag.

\begin{lstlisting}[caption={Motivating example: ALTER\_ERROR\_BRANCH},
  label=lst:motivating, basicstyle=\small\ttfamily,
  frame=single, numbers=left]
// Original JCL
//STEP02 EXEC PGM=CBACT02C,
//  COND=(4,GT,STEP01)   <- skip if RC > 4

// Modernized (mutation applied)
//STEP02 EXEC PGM=CBACT02C,
//  COND=(16,GT,STEP01)  <- threshold raised
\end{lstlisting}

%%=======================================================================
\section{Background}\label{sec:background}
%%=======================================================================

\textbf{JCL/COBOL batch systems.}
IBM Job Control Language (JCL) defines batch job structure: EXEC cards (program
invocations), DD cards (dataset declarations with DISP= disposition parameters),
and COND parameters (conditional execution based on step return codes).
COBOL programs implement business logic, referencing files via FILE CONTROL sections.
Behavioral obligations are distributed across artifact types and never written down
as specifications.

\textbf{Behavioral contracts.}
We extend Meyer's Design by Contract~\cite{meyer1992applying} to define a
\emph{behavioral contract} as a formally representable obligation that must hold
across execution events for the system to preserve its intended behavior.
Unlike DbC contracts, ours are \emph{recovered} from existing artifacts, not
specified by developers.

\textbf{Static vs.\ dynamic analysis.}
Static analysis extracts facts without execution; dynamic analysis extracts facts
from traces.
Both have known limitations: static over-approximates; dynamic under-approximates.
Our critical principle: $\notobs \neq \text{CONTRADICTED}$.
The absence of dynamic evidence for a contract is not evidence of absence; it signals
incomplete execution coverage.

\textbf{Static/dynamic reconciliation.}
Naik \etal~\cite{naik2006race} combine static and dynamic analysis to reduce
false positives in race detection.
BCG extends this pattern to behavioral contracts: static extraction produces candidate
contracts, dynamic reconciliation assigns explicit validation states
(\supported, \conflict, \notobs, \ambiguous, \partiallysupported) per contract, and
confidence-delta differencing drives modernization-risk classification.

%%=======================================================================
\section{Behavioral Contract Graph}\label{sec:bcg}
%%=======================================================================

\subsection{Formal Definition}

We define the \bcg as the 6-tuple:
$\mathit{BCG} = (V, E, C, P, D, \Gamma)$
where:
$V$ is the set of typed entities (jobs, steps, programs, datasets, copybooks);
$E \subseteq V \times \mathit{RelationType} \times V$ is the structural relation set;
$C$ is the behavioral contract set;
$P: C \to \mathit{ProvenanceRecord}$ maps each contract to its evidence record;
$D$ is the dynamic observation set (linked to $V$ via $\mathit{artifact\_id}$);
$\Gamma: C \to (\mathit{ValidationState} \times \mathit{ConfidenceVector})$ provides
per-contract validation state and evidence-strength scores.

\subsection{Contract Types}

\textbf{TEMPORAL:} $A \prec B$ — in all valid executions satisfying relevant
preconditions, $A$ must complete before $B$ begins.

\textbf{RESOURCE:} $\mathit{produces}(A, R) \prec \mathit{consumes}(B, R)$ — step
$A$ must successfully produce resource $R$ before step $B$ attempts to access it.

\textbf{CONTROL-FLOW:} $\mathit{RC}(A)\ op\ T \to \neg\mathit{execute}(B)$ — if
step $A$'s return code satisfies the condition $\mathit{RC}\ op\ T$, step $B$ must
not execute.

\subsection{Validation States and Confidence}

\Cref{table:states} defines the eight validation states.
The key scientific principle is that \notobs must remain distinct from
\texttt{CONTRADICTED}: a contract with no dynamic evidence is a valid static
candidate, not a falsified one.

The confidence model is a multi-dimensional evidence-strength vector
$\boldsymbol{w} \cdot \boldsymbol{x}$ with dimensions for static support, dynamic
support, cross-artifact support, provenance quality, human validation, and execution
frequency.
\textbf{Warning:} this score is an evidence-strength estimate in $[0,1]$, not a
calibrated probability.

\begin{table}[t]
\small\centering
\caption{Contract validation states.}\label{table:states}
\begin{tabular}{ll}
\toprule
State & Meaning \\
\midrule
\texttt{STATIC\_CANDIDATE}         & Static only; not yet dynamically validated \\
\texttt{STATIC\_DYNAMIC\_SUPPORTED} & Static and dynamic observations agree \\
\texttt{STATIC\_DYNAMIC\_CONFLICT}  & Static and dynamic observations disagree \\
\texttt{NOT\_OBSERVED}             & No dynamic evidence (≠ contradicted) \\
\texttt{HUMAN\_VALIDATED}          & Expert-confirmed \\
\texttt{GROUND\_TRUTH}             & Deterministic/execution-oracle verified \\
\texttt{DYNAMICALLY\_OBSERVED}     & Observed in traces; not statically extracted \\
\texttt{UNRESOLVED}                & Ambiguous; insufficient evidence \\
\bottomrule
\end{tabular}
\end{table}

%%=======================================================================
\section{Confidence-Delta Differencing}\label{sec:diff}
%%=======================================================================

The \texttt{BCGDiff} algorithm compares an original \bcg with a modernized \bcg.
It classifies each changed contract as either:
(a) \emph{broken}: the contract is absent from the modernized BCG; or
(b) \emph{modified}: the contract is present but its confidence has decreased
    below a threshold $\delta$.

The key insight is that a modernization engineer can \emph{weaken} a behavioral
obligation without removing it --- the \texttt{ALTER\_ERROR\_BRANCH} mutation in
\cref{lst:motivating} is precisely this case.
Structural analysis sees no change; BCG-without-confidence-delta (Baseline~E) also
misses it.
Only \texttt{diff.modified} inspection, checking
$\Delta\mathit{static\_support} = 1.0 - 0.7 = 0.3 > \theta$, raises the flag.

This architectural discriminator explains the full Baseline E $\to$ Proposed gap
(F1: $0.9600 \to 1.0000$), catching all 8 \texttt{ALTER\_ERROR\_BRANCH} mutations.

%%=======================================================================
\section{Static/Dynamic Reconciliation}\label{sec:reconciliation}
%%=======================================================================

For each contract $c$, the reconciliation engine retrieves events linked to
$c.\mathit{subject\_id}$ and $c.\mathit{object\_id}$ from $D$.
Predicate-specific logic determines the outcome:
\textsc{Precedes} compares latest-subject-timestamp with earliest-object-timestamp;
\textsc{ProducesBeforeConsumes} checks produce and consume event ordering;
\textsc{MustNotFollowOnError} checks whether the bypass condition is respected.

If no relevant events exist ($E_s = \emptyset \wedge E_o = \emptyset$):
outcome $\leftarrow$ \notobs (never \texttt{CONTRADICTED}).

The execution evidence must be labeled with its source type:
\texttt{EXECUTION} (real \gnucobol),
\texttt{EXECUTION\_SIMULATED} (JCL step-ordering simulation), or
\texttt{SYNTHETIC\_STATIC} (excluded by guard \texttt{SI-01}).
These labels must never be changed after assignment.

%%=======================================================================
\section{Fault-Grounded Mutation Framework}\label{sec:mutations}
%%=======================================================================

Mutation operators are derived from documented modernization failure modes, not tuned
to this benchmark.
We used the AWS \carddemo issue history, public migration case studies, and academic
literature~\cite{comella2000survey,ulrich2010information}.
The taxonomy was constructed before operator implementation and was not modified after
seeing evaluation results.

Three failure categories (F1--F3) with eight operators
(\cref{table:mutations}) cover job-ordering, file-availability, and
return-code-propagation regressions.
Of 440 mutation attempts, 336 (76.4\%) were not applicable
(\texttt{MutationNotApplicable}): F2/F3 operators require specific JCL structures
absent in most \carddemo files.
This is a benchmark representativeness limitation, not a framework failure.

%%=======================================================================
\section{Experimental Methodology}\label{sec:methodology}
%%=======================================================================

\textbf{Dataset.}
Primary: AWS \carddemo (Apache 2.0) --- 55 JCL files, 44 COBOL programs, 62
copybooks.
RQ1 uses 33 held-out JCL files (test split, seed=42).
RQ2 uses the development split.
Secondary: Open Mainframe Project public COBOL (43 JCL + 7 COBOL) for parser
coverage; not used as headline result.

\textbf{Ground truth hierarchy.}
Level-2 (deterministic structural facts from JCL grammar) for RQ2;
Level-1 (independent deterministic oracle) for RQ1 TEMPORAL contracts.
Level-3 (human annotation) is future work for RESOURCE/CONTROL\_FLOW.

\textbf{Baselines.}
\begin{description}
\item[A] Structural dependency analysis (program slicing)~\cite{tip1995slicing}.
\item[B] Static complexity/risk metrics.
\item[C] LLM-only: Groq \texttt{llama-3.3-70b-versatile}, v2\_blinded prompts,
  temperature=0.0.
\item[D] Dependency + LLM: topology provided to same LLM, v2\_blinded.
\item[E] Static \bcg without confidence-delta detection (\texttt{diff.broken} only).
\end{description}

\textbf{LLM baseline methodology.}
Baselines C and D used Groq's OpenAI-compatible API.
During validation, an initial prompt configuration (v1) was found to include the
phrase ``semantics-preserving change'' in negative descriptions, exposing benchmark
structure to the model.
Those runs were discarded.
All reported LLM results use v2\_blinded prompts with neutral operator descriptions.
High API error rates were observed: 245/248 (Baseline~C) and 241/248 (Baseline~D)
calls failed.
The reported F1 scores (0.038 and 0.109) are based on 3 and 7 successful calls
respectively, and should be interpreted as lower bounds.

\textbf{Negative benchmark.}
144 semantics-preserving mutations (5 operators: \texttt{RENAME\_STEP\_COMMENT},
\texttt{ADD\_REDUNDANT\_DEPENDENCY}, \texttt{REORDER\_INDEPENDENT\_STEPS},
\texttt{CHANGE\_ACCOUNTING\_PARAMETERS}, \texttt{RENAME\_PROC\_LABEL}) confirm
FPR=0.000.

\textbf{Dynamic evidence.}
40 execution traces: 2 real \gnucobol \texttt{EXECUTION}, 38
\texttt{EXECUTION\_SIMULATED} (JCL step-ordering simulation).
3 \texttt{SYNTHETIC\_STATIC} traces excluded by guard \texttt{SI-01}.
Execution environment: \gnucobol 3.2.0 on macOS arm64 --- \textbf{not} z/OS.

\textbf{Statistical method.}
95\% bootstrap CI (10,000 iterations, seed=42) for RQ1 macro-F1.
RQ2 metrics are descriptive (closed-set benchmark with FP=0 by design;
ceiling effects preclude pairwise significance tests).
Where significance testing is not appropriate, we say so explicitly.

\textbf{Engine fix disclosure.}
A gap in BCGDiff was identified and corrected before final evaluation: the original
\texttt{RiskDetector} inspected only \texttt{diff.broken}, missing mutations that
reduced contract confidence without removing the contract.
The fix added \texttt{diff.modified} inspection.
All RQ2 results were computed after this fix (\texttt{rq2\_summary.json},
\texttt{engine\_fix\_applied}).

\textbf{Reproducibility.}
\texttt{make reproduce-paper1} reproduces all deterministic results.
LLM baselines require \texttt{GROQ\_API\_KEY}.
All seeds fixed at 42.

%%=======================================================================
\section{Results}\label{sec:results}
%%=======================================================================

\subsection{RQ1: Contract Recovery}

\Cref{table:rq1} shows the TEMPORAL evaluation against the independently implemented
oracle (\texttt{scripts/run\_rq1\_independent.py}; no shared code with
\texttt{StaticContractExtractor}).

\textbf{Result:} Macro-F1 $= 0.576$ (95\% CI: [0.39, 0.73]);
micro-F1 $= 1.000$ (56 TP, 0 FP, 0 FN, TEMPORAL contracts only).
H1 is \textbf{supported}.

Macro-F1 is lower than micro-F1 because 17 of 33 held-out files have zero TEMPORAL
contracts; those files score F1=0.0 by convention.
This is not an extraction failure.

\textbf{Benchmark negative composition.}
The Level-2 ground-truth benchmark (463 items: 139 positive, 324 negative) contains
two negative classes: 113 \emph{reverse-temporal} negatives (hard: reversed step
order) and 208 \emph{self-PRECEDES} negatives (trivially false: a step before itself).
The self-PRECEDES negatives are detected trivially by any reflexivity guard and are
tagged \texttt{negative\_type=self\_precedes\_trivial} in the benchmark JSON.
RQ2 specificity (Specificity=1.000, 144 TN) uses a \emph{separate} mutation-based
negative benchmark containing only semantics-preserving mutations, not these GT negatives.
RQ1 evaluation uses the independent oracle only on 33 held-out files; no GT negatives
affect the primary RQ1 metric.

\textbf{Scope limitation:}
This evaluation covers TEMPORAL contracts only.
RESOURCE ($n=3$, single file) and CONTROL\_FLOW ($n=23$, no negatives)
are reported as counts only; their volumes are insufficient for statistical claims.
No independent oracle can be constructed for RESOURCE/CONTROL\_FLOW without sharing
parsing logic; Level-3 human annotation is deferred to future work.

\begin{table}[t]
\small
\centering
\caption{Contract recovery (RQ1) --- TEMPORAL contracts only, independently implemented oracle. The oracle (\texttt{scripts/run\_rq1\_independent.py}) shares only the JCL specification with the extractor; no source code is shared. Macro-F1 is lower because files with zero contracts contribute F1=0.0 by convention (17/33 such files). RESOURCE and CONTROL\_FLOW contracts require Level-3 human annotation for independent validation and are not reported here.}
\label{table:rq1}
\begin{tabular}{lrrrrrr}
\toprule
Scope & TP & FP & FN & Micro-P & Micro-R & Micro-F1 \\
\midrule
TEMPORAL (independent oracle) & 56 & 0 & 0 & 1.000 & 1.000 & 1.000 \\
\midrule
\multicolumn{7}{l}{\small Macro-F1 = 0.576 \quad 95\% CI (bootstrap, $n$=10{,}000): [0.39, 0.73] \quad $n_{\mathrm{files}}$=33} \\
\multicolumn{7}{l}{\small H1 verdict: \textbf{SUPPORTED} (threshold: Micro-F1 $\geq$ 0.70, Precision $\geq$ 0.75)} \\
\bottomrule
\end{tabular}
\vspace{0.3em}
{\small \textit{Scope limitation:} This evaluation is limited to TEMPORAL contracts. RESOURCE ($n$=3 in full corpus) and CONTROL\_FLOW ($n$=23) require independent human annotation; a deterministic oracle analogous to the above does not exist for those types without sharing parsing logic. The deprecated Level-2 self-consistency result (same extractor, F1=1.00) is preserved in \texttt{results/rq1/rq1\_summary.json} but is not cited as accuracy evidence.}
\end{table}


\subsection{RQ2: Modernization Risk Detection}

\Cref{table:baselines,table:per_operator} present the full results.

\textbf{Positive benchmark (104 mutations; F2 operators n=1 each — preliminary):}
Static \bcg achieves F1=1.000 (recall=1.000, FNR=0.000).
The discriminating finding is \texttt{ALTER\_ERROR\_BRANCH}: 8 mutations where
confidence-delta detection identifies a modified (not broken) contract with
$\Delta\mathit{static\_support} = 0.30$.
Structural Baseline A (F1=0.900) and BCG-without-confidence-delta Baseline E
(F1=0.960) both miss these; the proposed method catches all 8.
F2 operators (REMOVE\_DATASET\_INIT, CHANGE\_DATASET\_NAME,
ALTER\_FILE\_DEP\_ORDER) each have $n=1$ applicable instance;
those results are illustrative and not sufficient for per-operator statistical claims.

\textbf{Negative benchmark (144 TN):}
FPR=0.000, Specificity=1.000 for all methods.
\textit{Note: FPR=0.000 is a consequence of the closed-set benchmark design;
open-world false-positive rates are not measured.}

\textbf{LLM baselines:}
On this blinded benchmark, Baseline~C (LLM-only) achieved F1=0.038 and Baseline~D
(Dependency+LLM) achieved F1=0.109.
These results are lower bounds due to high API error rates; they should not be
generalized beyond this benchmark, task, and model version.

H2 (reformulated) is \textbf{supported}:
recall(proposed) $= 1.00 >$ recall(Baseline A) $= 0.82$ and recall(proposed) $= 1.00 >$
recall(Baseline E) $= 0.92$.

\begin{table}[t]
\small
\centering
\caption{Baseline comparison (RQ2). Evaluated on 104 applicable mutations (440 attempted, 336~N/A). FP=0 by construction. FPR from negative benchmark: TN=144, FP=0, FPR=0.000. LLM baselines (C, D) use Groq \texttt{llama-3.3-70b-versatile}, v2\_blinded prompts (v1 discarded due to benchmark leak); 245/241 API errors reduced effective sample. All non-LLM results are deterministic and fully reproducible.}
\label{table:baselines}
\begin{tabular}{lrrrrrrl}
\toprule
Method & TP & FP & FN & P & R & F1 & FNR \\
\midrule
Complexity (Baseline B)       & 0  & 0 & 104 & 0.000 & 0.000 & 0.000 & 1.000 \\
Structural (Baseline A)       & 85 & 0 &  19 & 1.000 & 0.817 & 0.900 & 0.183 \\
Static BCG w/o $\delta$ (Baseline E) & 96 & 0 & 8 & 1.000 & 0.923 & 0.960 & 0.077 \\
LLM-only (Baseline C)$^\dagger$ & 2  & 0 & 102 & 1.000 & 0.019 & 0.038 & 0.981 \\
Dep+LLM (Baseline D)$^\dagger$ & 6  & 0 &  98 & 1.000 & 0.058 & 0.109 & 0.942 \\
\midrule
\textbf{Static BCG (Proposed)} & \textbf{104} & \textbf{0} & \textbf{0} & \textbf{1.000} & \textbf{1.000} & \textbf{1.000} & \textbf{0.000} \\
\bottomrule
\end{tabular}
\vspace{0.3em}
{\small $^\dagger$High API error rates (245/248 and 241/248 calls failed); classification based on 3 and 7 successful responses, respectively. Low recall driven primarily by API failures, not by misclassification. This limitation is discussed in \S\ref{sec:threats}.}
\vspace{0.3em}
{\small Negative benchmark (144 semantics-preserving mutations): FP=0 for all methods evaluated; FPR=0.000; Specificity=1.000.}
\end{table}


\subsection{RQ3: Dynamic Evidence}

\Cref{table:rq3} shows the reconciliation results with 40 execution traces.

\textbf{Result:}
56/61 contracts reach \supported (91.8\%).
1 contract reaches \conflict (a CONTROL\_FLOW contract where trace shows execution
despite COND-bypass condition; likely GnuCOBOL/z/OS behavioral divergence).
4 contracts remain \notobs (JCL steps requiring DB2/CICS execution environment).

\textbf{Coverage experiment (\cref{table:coverage}):}
Validation rate increases monotonically from 1.6\% (4 traces) to 91.8\% (40 traces).
No \notobs contract becomes \texttt{CONTRADICTED} as traces are withheld
(\texttt{not\_observed\_monotonicity\_safe = True}).
This empirically confirms $\notobs \neq \text{CONTRADICTED}$.

H3 (reformulated) is \textbf{supported}: dynamic validation rate $= 91.8\% > 50\%$;
FPR $= 0.000$ (TN=144).

\begin{table}[t]
\small\centering
\caption{RQ3: Dynamic contract reconciliation (40 traces).}\label{table:rq3}
\begin{tabular}{lrr}
\toprule
State & Count & Rate \\
\midrule
\texttt{STATIC\_DYNAMIC\_SUPPORTED} & 56 & 91.8\% \\
\texttt{STATIC\_DYNAMIC\_CONFLICT}  & 1  & 1.6\%  \\
\texttt{NOT\_OBSERVED}              & 4  & 6.6\%  \\
\midrule
Total contracts & 61 & — \\
\bottomrule
\end{tabular}
\end{table}

\begin{table}[t]
\small\centering
\caption{RQ3: Dynamic coverage experiment.}\label{table:coverage}
\begin{tabular}{rrrrrr}
\toprule
Coverage & Traces & Supported & Conflict & NOT\_OBS & Rate \\
\midrule
10\%  & 4  & 1  & 0 & 60 & 1.6\%  \\
25\%  & 10 & 15 & 0 & 46 & 24.6\% \\
50\%  & 20 & 26 & 0 & 35 & 42.6\% \\
75\%  & 30 & 39 & 0 & 22 & 63.9\% \\
100\% & 40 & 56 & 1 & 4  & 91.8\% \\
\bottomrule
\end{tabular}
{\small NOT\_OBS $\neq$ CONTRADICTED at all levels (monotonicity confirmed).}
\end{table}

%%=======================================================================
\section{Discussion}\label{sec:discussion}
%%=======================================================================

\textbf{Why does BCG outperform structural analysis?}
Structural analysis identifies topology changes but cannot represent behavioral
obligations.
The \texttt{ALTER\_ERROR\_BRANCH} family demonstrates this precisely: topology is
unchanged, but the behavioral obligation (``STEP02 must not execute if STEP01 returns
8--15'') is violated.
\bcg represents this as a \textsc{Control-Flow} contract with a confidence that
degrades under the mutation; structural analysis has no analogous representation.

\textbf{Why does confidence-delta matter?}
Modernization risks are not always binary.
A practitioner may weaken a guard condition (raise the RC threshold) rather than
remove it entirely.
The \texttt{diff.modified} mechanism detects this weakening class, which accounts for
8 mutations (7.7\% of evaluated cases) that Baseline E misses.

\textbf{Why did the LLM baselines perform poorly on this benchmark?}
On this blinded benchmark, F1 $= 0.038$ (LLM-only) and F1 $= 0.109$ (Dep+LLM).
Both results are lower bounds due to high API error rates (97--98\%).
At the same time, even the 3 and 7 successful calls produced modest recall.
The blinded baselines lack access to the BCG's formal contract representation,
confidence model, and structured diff infrastructure.
We do not claim that LLMs generally fail on semantic analysis tasks; these results
are specific to this benchmark, task formulation, and model/prompt configuration.

\textbf{What does dynamic evidence add?}
Dynamic evidence promotes 91.8\% of statically extracted TEMPORAL contracts to
\supported, confirming that JCL step-ordering obligations are valid execution
obligations, not artifacts of static analysis.
The 1 \conflict case warrants investigation; the 4 \notobs cases are definitionally
not contradictions.
The coverage experiment demonstrates that each additional trace provides new evidence
without creating false contradictions.

\textbf{Deployment on real z/OS.}
Production deployment would require: (1) z/OS SMF-record-based trace collection;
(2) a JES job-completion event adapter; (3) a RACF-aware entity resolver.
None of these are fundamental architectural changes; they are adapter implementations
for the existing \texttt{DynamicEvidenceAdapter} interface.

%%=======================================================================
\section{Threats to Validity}\label{sec:threats}
%%=======================================================================

\textbf{Internal validity.}
Mutation operators were designed from external evidence before BCG implementation;
no operator was modified after seeing results.
However, the benchmark is dominated by F1 (TEMPORAL) mutations (85/104); F2/F3
coverage is limited to 19 cases.
The LLM baseline results are lower bounds due to API infrastructure constraints.

\textbf{Construct validity.}
The independent oracle for RQ1 covers TEMPORAL contracts only; RESOURCE and
CONTROL\_FLOW require human annotation.
The confidence score is an evidence-strength estimate, not a calibrated probability.
FP=0 in RQ2 is a closed-benchmark design property, not an independently measured
false-alarm rate.

\textbf{External validity.}
All results are on \carddemo --- a demonstration system of modest complexity
($\approx30$ COBOL programs, 33+ JCL files).
Real enterprise systems may have thousands of programs and more complex JCL.
We do not claim enterprise-wide generalization.
A second public JCL/COBOL system was not empirically evaluated; this remains an
explicit limitation.

\textbf{Out-of-distribution generalization (E5).}
We tested the BCG extractor on a public COBOL corpus (43 JCL files, public GitHub
repositories) as an out-of-distribution (OOD) probe.
The in-distribution F1 on \carddemo is $1.000$; the OOD F1 is $0.000$
(precision=$0.000$, 34 false positives, 0 true positives).
Root-cause analysis shows the OOD corpus lacks the JCL orchestration structure
required for temporal contracts: the extractor correctly identifies temporal
ordering in the JCL files present, but those contracts have no counterpart in the
OOD ground truth (which was constructed for a COBOL-only, non-orchestrated corpus).
This result does not indicate an extractor defect; it indicates that BCG is
applicable to systems with \emph{joint} JCL+COBOL artifact structure and should
not be applied to COBOL-only corpora without JCL orchestration context.
Scope: all claims are restricted to systems exhibiting JCL job-stream orchestration.

\textbf{GnuCOBOL vs.\ z/OS.}
38/40 traces are \texttt{EXECUTION\_SIMULATED}; 2 are real \gnucobol 3.2.0.
z/OS-specific behaviors (SMF, JES, RACF, COMP-3 binary arithmetic) are not captured.
All dynamic evidence is labeled \texttt{EXECUTION (GnuCOBOL, NOT z/OS)}.

\textbf{LLM validity.}
LLM results depend on model version, prompt, and provider.
Infrastructure constraints (97--98\% API error rate) make the reported F1 values
lower bounds.
These results apply only to this benchmark, task, and configuration.

%%=======================================================================
\section{Related Work}\label{sec:related}
%%=======================================================================

\textbf{Dynamic invariant inference.}
Daikon~\cite{ernst2007daikon} infers data-value invariants from execution traces.
BCG differs on four axes: (a) contract type (behavioral obligations vs.\ data-value
invariants); (b) artifact domain (cross-artifact JCL+COBOL vs.\ single-program OO);
(c) static extraction precedes dynamic grounding; (d) formal provenance per contract.

\textbf{Specification mining.}
Ammons \etal~\cite{ammons2002mining} mine API usage protocols from program traces;
Lo \etal~\cite{lo2006mining} mine iterative-convergent control patterns;
Lorenzoli \etal~\cite{lorenzoli2008automatic} (\textsc{Gist}) infer behavioral
models via FSM learning from execution sequences;
Klaedtke \etal~\cite{klaedtke2021mining} mine behavioral contracts from execution
traces using DFA inference;
Mariani and Pezz\`{e}~\cite{mariani2012behavioural} extend this to web applications.
Angluin~\cite{angluin1987learning} provides the foundational $L^*$ algorithm
for regular-set inference underpinning these approaches.
BCG extends this lineage to cross-artifact legacy batch systems, adds static
evidence extraction that precedes dynamic reconciliation, introduces an explicit
conflict taxonomy, and provides formal provenance per contract.

\textbf{\textsc{Gist} vs.\ BCG.}
Lorenzoli \etal's \textsc{Gist}~\cite{lorenzoli2008automatic} is the closest
dynamic contract-inference predecessor. BCG differs on five axes:
(1)~\emph{Static+dynamic reconciliation:} \textsc{Gist} is purely dynamic; BCG
extracts static structural evidence first, then reconciles with traces via explicit
conflict states;
(2)~\emph{Provenance:} \textsc{Gist} FSMs carry no link to source artifacts; BCG
contracts carry a \texttt{ProvenanceRecord} per contract;
(3)~\emph{Cross-artifact scope:} \textsc{Gist} targets single programs; BCG spans JCL
jobs, COBOL steps, and inter-artifact dataset dependencies;
(4)~\emph{Batch domain semantics:} \textsc{Gist} models method-call event traces; BCG
models JCL step ordering, DD \texttt{DISP=} parameters, and \texttt{COND=} return-code
chains;
(5)~\emph{Fault-grounded evaluation:} \textsc{Gist} was not evaluated with
empirically sourced modernization mutation operators; BCG operators derive from
documented mainframe-to-cloud migration failures.
Similarly, Klaedtke \etal's DFA-based mining~\cite{klaedtke2021mining} is dynamic-only,
lacks static provenance and cross-artifact scope, and targets general program
traces rather than legacy batch orchestration.

\textbf{Runtime verification.}
Leucker and Schallhart~\cite{leucker2009brief} survey runtime verification as
inline specification checking.
BCG performs \emph{post-hoc} recovery for modernization risk analysis, not
inline monitoring.

\textbf{Semantic differencing.}
Lahiri \etal~\cite{lahiri2012regression} (\textsc{Symdiff}) perform statement-level
symbolic differencing;
Felsing \etal~\cite{felsing2014automating} automate regression verification via
deductive self-composition.
BCG operates at job/step/contract granularity without SMT encoding (infeasible for
large COBOL programs with vendor runtime behavior).
Pessoa \etal~\cite{pessoa2022behavioral} perform behavioral semantic differencing on
COBOL program pairs to detect modernization regressions.
Unlike Pessoa, BCG: (a) represents behavioral obligations as a formal typed contract IR
with per-contract provenance and validation states; (b) spans JCL job-level orchestration
and batch step ordering across multiple programs; (c) reconciles static structural
evidence with dynamic execution traces; (d) evaluates on a fault-grounded mutation
benchmark derived from documented production failures.

\textbf{Legacy COBOL analysis.}
Sneed~\cite{sneed2008cobol} extracts structural facts from COBOL for migration
assessment; our work adds behavioral obligations, JCL scope, and dynamic evidence.
Comella-Dorda \etal~\cite{comella2000survey} and Ulrich and Newcomb~\cite{ulrich2010information}
document the failure patterns our mutation taxonomy is grounded in.

\textbf{LLM code analysis.}
Hou \etal~\cite{hou2024llm4se}, Fan \etal~\cite{fan2023llm}, and Ling \etal~\cite{ling2023domain}
demonstrate LLM code-understanding capabilities.
Our evaluation provides a datapoint on LLM performance for blinded behavioral
contract violation detection in legacy JCL/COBOL.

%%=======================================================================
\section{Limitations}\label{sec:limitations}
%%=======================================================================

\begin{description}
\item[L1] \textbf{Single primary dataset.} All results are on \carddemo.
  Enterprise systems may have thousands of programs; performance may differ.
\item[L2] \textbf{GnuCOBOL, not z/OS.} Dynamic evidence is \gnucobol on
  Linux/macOS; z/OS-specific behaviors are not captured.
\item[L3] \textbf{Simulated traces.} 38/40 traces are \texttt{EXECUTION\_SIMULATED};
  only 2 are real execution.
\item[L4] \textbf{TEMPORAL-only independent validation.} RESOURCE and CONTROL\_FLOW
  contracts require human annotation; oracle scope is explicitly limited.
\item[L5] \textbf{LLM results are benchmark-specific.}
  F1=0.038/0.109 applies to this benchmark, task, model, and prompt; high API error
  rates make these lower bounds.
\item[L6] \textbf{Confidence is not probability.} The confidence score is an
  evidence-strength estimate without calibration.
\end{description}

%%=======================================================================
\section{Conclusion}\label{sec:conclusion}
%%=======================================================================

We presented the Behavioral Contract Graph (\bcg), a provenance-aware intermediate
representation that makes implicit behavioral obligations explicit in legacy
COBOL/JCL systems, and confidence-delta differencing (\texttt{BCGDiff}) for
modernization risk detection.

Our empirical evaluation on the AWS \carddemo benchmark produced three findings:
(1) The TEMPORAL extraction pipeline achieves micro-F1=1.000 against an independent
oracle (56 TP, 0 FP, 0 FN; macro-F1=0.576, CI:[0.39,0.73]);
(2) Static \bcg detects all 104 evaluated positive mutations (F1=1.000) with
FPR=0.000 on 144 semantics-preserving cases; the key contributor is confidence-delta
detection of 8 \texttt{ALTER\_ERROR\_BRANCH} mutations invisible to structural and
BCG-only baselines;
(3) Dynamic evidence validates 91.8\% of contracts, and the coverage experiment
confirms $\notobs \neq \text{CONTRADICTED}$ empirically.

Limitations include single-dataset evaluation, GnuCOBOL rather than z/OS execution,
and TEMPORAL-only independent oracle coverage.
These are explicitly retained, not minimized.

%%=======================================================================
\section*{Data Availability}\label{sec:data}
%%=======================================================================

All experimental results, scripts, and configurations are available in the
supplementary artifact package.
Deterministic components are fully reproducible via \texttt{make reproduce-paper1}.
LLM baselines require a \texttt{GROQ\_API\_KEY} (not included).
The primary dataset is AWS \carddemo (Apache 2.0, publicly available).

%%=======================================================================
%% Bibliography
%%=======================================================================
\bibliographystyle{ACM-Reference-Format}
\bibliography{references}

%% -----------------------------------------------------------------------
%% Appendix: Per-Operator Results
%% -----------------------------------------------------------------------
\appendix
\section{Per-Operator Breakdown}\label{app:per_operator}

\begin{table}[t]
\small
\centering
\caption{Per-operator detection (RQ2). Bold = operators where structural baseline scores zero, demonstrating the need for behavioral contract semantics.}
\label{table:per_operator}
\begin{tabular}{llrrrr}
\toprule
Operator & Category & BCG F1 & Struct. F1 & $\Delta$ & N Eval. \\
\midrule
REORDER\_JOBS & F1 & 1.000 & 0.966 & +0.035 & 29 \\
REMOVE\_JOB\_DEPENDENCY & F1 & 1.000 & 1.000 & 0.000 & 28 \\
ALTER\_STEP\_SEQUENCE & F1 & 1.000 & 1.000 & 0.000 & 28 \\
REMOVE\_DATASET\_INIT & F2 & 1.000 & 1.000 & 0.000 & 1 \\
\textbf{CHANGE\_DATASET\_NAME} & F2 & \textbf{1.000} & \textbf{0.000} & \textbf{+1.000} & 1 \\
\textbf{ALTER\_FILE\_DEP\_ORDER} & F2 & \textbf{1.000} & \textbf{0.000} & \textbf{+1.000} & 1 \\
\textbf{SUPPRESS\_RC\_CHECK} & F3 & \textbf{1.000} & \textbf{0.000} & \textbf{+1.000} & 8 \\
\textbf{ALTER\_ERROR\_BRANCH} & F3 & \textbf{1.000} & \textbf{0.000} & \textbf{+1.000} & 8 \\
\midrule
\textbf{Overall} & & \textbf{1.000} & \textbf{0.817} & \textbf{+0.183} & \textbf{104} \\
\bottomrule
\end{tabular}
\end{table}


\section{Fault-Grounded Mutation Taxonomy}\label{app:mutations}

\begin{table}[t]
\small
\centering
\caption{Fault-grounded mutation taxonomy. All operators are empirically documented; see \texttt{docs/failure\_taxonomy.md}.}
\label{table:mutations}
\begin{tabular}{llrr}
\toprule
Operator & Failure Category & N Evaluated & N N/A \\
\midrule
REORDER\_JOBS & F1 — Job ordering & 29 & 26 \\
REMOVE\_JOB\_DEPENDENCY & F1 — Job ordering & 28 & 27 \\
ALTER\_STEP\_SEQUENCE & F1 — Job ordering & 28 & 27 \\
REMOVE\_DATASET\_INIT & F2 — Resource & 1 & 54 \\
CHANGE\_DATASET\_NAME & F2 — Resource & 1 & 54 \\
ALTER\_FILE\_DEP\_ORDER & F2 — Resource & 1 & 54 \\
SUPPRESS\_RC\_CHECK & F3 — RC propagation & 8 & 47 \\
ALTER\_ERROR\_BRANCH & F3 — RC propagation & 8 & 47 \\
\midrule
\textbf{Total} & & \textbf{104} & \textbf{336} \\
\bottomrule
\end{tabular}
\end{table}


\end{document}
